\documentclass[11pt,a4paper]{article}
\usepackage[margin=1in]{geometry}
\usepackage[T1]{fontenc}
\usepackage{titlesec}
\usepackage{enumitem}
\usepackage{xcolor}
\usepackage{fancyhdr}
\usepackage{hyperref}
\usepackage{parskip}

\hypersetup{
  colorlinks=true,
  linkcolor=blue!60!black,
  urlcolor=blue!60!black,
  citecolor=blue!60!black
}

\pagestyle{fancy}
\fancyhf{}
\fancyhead[L]{PetSocial --- User Guide}
\fancyfoot[C]{\thepage}
\renewcommand{\headrulewidth}{0.4pt}

\title{PetSocial\\\large User Guide}
\author{PetSocial}
\date{\today}

\begin{document}
\maketitle
\thispagestyle{fancy}
\tableofcontents

\section{Getting Started}

\subsection{Creating an account}
Open the app and choose \textbf{Sign up}. You will need a display name, a username
(this becomes part of your profile's web address), an email address, and a password.
Depending on how the app is configured, you may need to click a confirmation link sent
to your email before you can log in for the first time.

\subsection{Adding your first pet}
From your home screen, choose \textbf{Add a pet}. A pet profile is separate from your
own: give it a name and a species (for example ``Dog'' or ``Cat''); breed and a short
bio are optional. Once created, you can post as yourself \emph{or} as this pet.

\section{Posting and the Home Feed}
Tap \textbf{New post} from the navigation bar. Add a caption, a photo, or both, and
optionally choose to post as one of your pets instead of yourself. Your \textbf{Home}
feed shows posts from you and everyone you follow; if you are not signed in, or not
following anyone yet, you will see a general feed of recent posts instead. You can like
and comment on any post, and tapping a post opens its full detail with all comments.

\section{Following People}
Visit someone's profile --- your own pet or theirs --- and tap \textbf{Follow}. The
\textbf{Explore} tab lists other owners you can search for and follow, and includes
quick shortcuts to Communities, Lost \& Found, Book a Vet, and the Marketplace.

\section{Your Profile and Privacy}
\subsection{Editing your profile}
From your profile, choose \textbf{Edit profile} to change your display name, bio, or
photo.

\subsection{Staying anonymous}
If you would rather not show your real name and photo, turn on \textbf{Keep my
identity anonymous} in your profile settings. Your name will show as ``Anonymous''
everywhere in the app, on every post and pet you're linked to. You can also apply this
setting to a single pet only (in that pet's own edit screen) if you want that one pet's
posts to stay anonymous while you post under your real name elsewhere. Turning this on
never affects your ability to be messaged or to use the app --- it only changes what
other people see.

\section{Communities}
Every pet's species automatically becomes a community. Open \textbf{Communities} to
browse by species, or narrow further by breed once inside a species page, to see other
pets and recent posts from that group.

\section{Lost \& Found}
If your pet goes missing --- or you have found someone else's pet --- open
\textbf{Lost \& Found} and choose \textbf{Report a pet}. Provide a description and, if
you can, contact information. Once the situation is resolved, open your own report and
mark it resolved so it no longer appears in the active list.

\section{Health Log}
Open one of your own pets' profiles to see its \textbf{Health log}, visible only to
you. Use \textbf{Add entry} to record a vaccination, a vet visit, a medication, or
anything else worth tracking, with an optional reminder date.

\section{Notifications}
The bell icon shows you likes, comments, and new followers, with an unread count badge.
Use \textbf{Mark all read} to clear it.

\section{Messages}
Tap the message icon (or \textbf{Message} on someone's profile) to start a private
conversation. Your conversation list shows your most recent message with each person
and how many are unread; opening a conversation marks it read.

\section{Adopting a Pet}
Open \textbf{Adopt} to browse pets currently available, filterable by species. Tap a
listing to see its full description and contact the poster directly through
\textbf{Message about \emph{name}}. If you are the one posting an adoptable pet, use
\textbf{Post a pet} and update its status (Available / Pending / Adopted) from its
detail page as things progress.

\section{Booking a Vet Appointment}
Open \textbf{Book a Vet} to see a directory of clinics. Choose one and \textbf{Book}:
pick which of your pets the visit is for, a reason, a day within the next week, and one
of the times still open that day. Your appointments --- and a \textbf{Cancel} option
for any of them --- are listed under \textbf{My appointments}.

\section{Marketplace}
Open \textbf{Marketplace} to browse pet products for sale, filterable by category.
Selling something yourself uses \textbf{Sell something}: add a title, category, price,
description, and an optional photo. Buyers reach you the same way as elsewhere in the
app, through \textbf{Message about \emph{item}}; once something sells, mark the
listing \textbf{Sold} from its detail page --- it stays visible with a Sold label
rather than disappearing, so your listing history remains browsable.

\section{Getting Help}
If something in the app is not behaving as described here, or you believe you've found
an error, please report it through the support channel your app deployment provides.

\end{document}
